% ICASSP 2027 manuscript. Compile RCS.tex with pdfLaTeX + BibTeX.
\documentclass[letterpaper]{article}
\usepackage{spconf}
\usepackage{amsmath,amssymb,graphicx}
\usepackage{booktabs}
\usepackage{cite}
\usepackage{url}
\usepackage{iftex}
\ifPDFTeX
  \usepackage[utf8]{inputenc}
  \usepackage[T5,T1]{fontenc}
  \newcommand{\vi}[1]{{\fontencoding{T5}\selectfont #1}}
\else
  \usepackage{fontspec}
  \setmainfont{texgyretermes-regular.otf}[
    BoldFont=texgyretermes-bold.otf,
    ItalicFont=texgyretermes-italic.otf,
    BoldItalicFont=texgyretermes-bolditalic.otf]
  \newcommand{\vi}[1]{#1}
\fi
\usepackage{tikz}
\usetikzlibrary{arrows.meta,positioning}
\usepackage{pgfplots}
\pgfplotsset{compat=1.18}
\usepackage[hidelinks]{hyperref}

\newcommand{\method}{RCS}
\newcommand{\TR}{TR}
\newcommand{\rep}{\ensuremath{\mathrm{Rep}}}
\newcommand{\ins}{\ensuremath{\mathrm{Ins}}}
\DeclareMathOperator*{\argmax}{arg\,max}
\newcommand{\cE}{\mathcal{E}}
\newcommand{\cC}{\mathcal{C}}
\newcommand{\cV}{\mathcal{V}}
\newcommand{\cA}{\mathcal{A}}
\newcommand{\cD}{\mathcal{D}}
\newcommand{\cS}{\mathcal{S}}

\title{As Speakers Would Switch:\\
In-Place Editing and Native-Accent Respelling for Code-Switched ASR}
\name{Tien Dat Bui\textsuperscript{1},
Van Dat Le\textsuperscript{1},
Ngoc Thuan Tran\textsuperscript{1},
Nhat Minh Le\textsuperscript{1},
Van Hai Do\textsuperscript{2}\sthanks{Corresponding author}}
\address{\textsuperscript{1}Viettel AI, Viettel Group, Vietnam\\
\textsuperscript{2}Thuyloi University, Hanoi, Vietnam}

\begin{document}
\maketitle

\begin{abstract}
Automatic speech recognition remains weak on code-switched speech, largely
because natural code-switched data are scarce and existing text-generation
methods are not natural enough. Text-to-speech augmentation is cheaper, but
is bounded by the text it renders: canonical synthesis cannot reproduce the
diverse, matrix-language-adapted pronunciations that code-switching exhibits
in practice. We present RCS, which takes only monolingual audio and a set of
code-switching terms and generates natural, logically coherent code-switched
data for a new domain while reusing the colloquial style of the
matrix-language transcripts; a companion spelling augmentation models the
pronunciation variants of switched terms. Applied to Vietnamese--English,
RCS lowers WER by 22\% relative on both Whisper and Conformer-CTC and raises
term recall by 55\% in domain, with gains carrying to out-of-domain terms
absent from the catalog (+50\% on Whisper and +138\% on Conformer-CTC).
Its synthetic text is the closest to real code-switched speech among the
compared generators, with a masked PPL gap of 0.07 versus 0.17 and 0.68.
\end{abstract}
\begin{keywords}
code-switching, speech recognition, data augmentation, text-to-speech,
large language models
\end{keywords}

\section{Introduction}
\label{sec:intro}
Code-switching (CS) remains difficult for automatic speech recognition
(ASR), especially when lexically sparse names, products, and technical
terms are inserted into matrix-language
speech~\cite{zhou2026sorry,liu2026lens}. These terms are often the most
informative words in an utterance, so errors on them are
disproportionately costly~\cite{ugan2025pier}. In deployment the terms
that matter are largely known in advance, but the list changes over time,
while natural CS corpora remain limited and
domain-specific~\cite{nguyen2026vimedcss}; recording new speech for every
new term is impractical, making augmentation attractive.

Text-side augmentation generates CS text and synthesizes it with
TTS~\cite{yan2025csfleurs}, evolving from rule-based substitution under
linguistic constraints~\cite{pratapa2018} and GAN-based
generation~\cite{chang2019gan} to LLM
prompting~\cite{hu2023llmtext,singh2025leveraging,alharbi2024leveraging}
and constrained decoding~\cite{kuwanto2024,liu2026lens}; contextual
biasing~\cite{pundak2018deep,sun2021tree} instead injects known terms at
decoding time and is complementary to it. Two limitations remain. First,
full generation or rewriting replaces the surrounding conversational
register and lets the LLM choose an uncontrolled embedded vocabulary, and
its written-style output loses colloquial speech patterns. Second,
canonical embedded-language TTS
ignores the matrix-language-adapted pronunciations that bilingual
speakers routinely produce: embedded items undergo L1-conditioned
phonological transfer and deviate from their dictionary form.
Spelling augmentation is a natural
remedy, encoding these variants as phonetically plausible alternative
spellings, and prior work shows that perturbing reference transcriptions
with similar-sounding spellings improves rare-word
robustness~\cite{huang2024perturb} and that LLM-driven phonetic respelling
broadens lexical and phonetic coverage~\cite{yamashita2026synthetic}.

We propose \method{} (\textbf{R}ealistic \textbf{C}ontext
\textbf{S}ynthesis), instantiated here with Vietnamese as the matrix
language and English as the embedded language: it preserves each real
Vietnamese transcript and edits only selected switch positions. We make
two contributions: (1) \emph{catalog-grounded text variants}, where an
LLM proposes plausible replacement/insertion slots grounded to a curated
term catalog, so a single transcript yields many controlled CS variants
without asking the LLM to invent the English term, and new terms can be
added without new speech or new slot proposals; and (2) \emph{respelling
augmentation}, where embedded terms are stochastically rendered with
Vietnamese-adapted spellings before TTS while training labels remain
canonical.

\begin{figure*}[t]
    \centering
    \includegraphics[width=\textwidth]{rcs_pipeline_v2 (1).pdf}
    \caption{RCS separates switch-location proposal from catalog-based
    term filling. Respelling changes the TTS input $z$, while the ASR
    label $y$ remains canonical. All text outside the selected edits is
    preserved; the source recording supplies the speaker reference.}
    \label{fig:pipeline}
\end{figure*}

\section{Method}
\label{sec:method}

\method{} edits real transcripts in two stages (catalog-grounded slot editing and optional respelling) before synthesis (Fig.~\ref{fig:pipeline}).

\subsection{Problem formulation}
Let $\cD_{\rm mono}=\{(x_i,T_i)\}_{i=1}^{N}$ be a monolingual matrix-language corpus with transcripts $T=(w_1,\ldots,w_n)$. A term catalog $\cE=\{(c,d_c,\cV_c)\}$ associates each category $c$ with a one-line description $d_c$ and a set of values $\cV_c$ (e.g., $\cV_{\rm tech}=\{$iphone, macbook, server$\}$). \method{} is a stochastic map $T\mapsto(y,z)$ that produces a CS label $y$ and a synthesis string $z$, yielding a synthetic pair $(\tilde x,y)$ with $\tilde x=\mathrm{TTS}(z;\,x)$, where the source recording $x$ serves as speaker reference. The label is the source transcript with at most $K$ localized edits applied in place,
\begin{equation}
  y = T \oplus \{(a_k,e_k)\}_{k=1}^{K}, \qquad e_k\in\textstyle\bigcup_{c}\cV_c,
  \label{eq:edit}
\end{equation}
where $a_k$ is either a span to be replaced (\rep) or an insertion anchor (\ins), and $\oplus$ applies the non-overlapping edits. Everything outside the $K$ edited positions, including disfluencies and discourse markers, is copied from the real transcript. An ASR model is finally trained on $\cD_{\rm mono}$ together with the synthetic set $\cD_{\rm syn}$.



\subsection{LLM-guided slot editing}
Given $T$, an LLM returns a set of proposals $P(T)=\{p_j\}$ with $p_j=(o_j,a_j,t_j,\delta_j)$: an operation $o_j\in\{\rep,\ins\}$, a span or anchor $a_j$, a free-form semantic tag $t_j$, and a short description $\delta_j$.We use Gemini~2.5 Flash with a chain-of-thought prompt: the model first
identifies concrete, nameable spans and assigns each an entity type, preferring \rep{} and falling back to a single-word \ins{} only when no replaceable span exists. It then rejects proposals that change the
grammatical role of the span, alter the meaning, or target abstract emotion and state words that bilingual speakers essentially never switch, before emitting a JSON list of proposals. The LLM never writes the CS sentence and never chooses the English term; it contributes span selection and grammatical judgment only.

\subsection{Catalog grounding and variant expansion}
Free-form tags are unstable across calls (\textsc{SportsClub}, \textsc{FootballTeam}, and \textsc{TeamName} denote one class), and the terms the LLM volunteers are not always real. Each proposal is therefore grounded in $\cE$ by semantic retrieval,
\begin{equation}
  \hat c_j=\argmax_{c\in\cE}\cos\!\big(\phi(t_j\Vert\delta_j),\,\phi(c\Vert d_c\Vert\cV_c)\big),
  \label{eq:ground}
\end{equation}
where $\phi$ is a general-purpose text embedding model (gemini-embedding-001) and $\Vert$ denotes concatenation. A slot is admitted only if the maximum similarity exceeds a threshold $\tau$; no category is ever created on the fly. Denoting the admitted slots by $\cA(T)$, the variant set is
\begin{equation}
\begin{aligned}
\mathcal{V}(T)
= \{\,T \oplus \{(a_j,e_j)\}_{j\in S}:\;&
S \subseteq A(T),\ |S|\le K,\\
& e_j\in V_{\hat c_j}\,\},
\end{aligned}
\label{eq:variant}
\end{equation}
over non-overlapping subsets $S$, so that
\[
|\mathcal{V}(T)|
=
\sum_{S} \prod_{j\in S} |V_{\hat c_j}|,
\]
where the sum is over all non-overlapping $S\subseteq A(T)$ with $|S|\le K$.
Diversity thus accumulates along three axes (alternative values
within a category, alternative switch locations, and both \textsc{Rep} and \textsc{Ins}), and it is bounded by the richness of the catalog rather than by the generator. We draw $M$ variants per transcript without replacement. Because admitted slots and their categories are cached, a new term added to $\cV_c$ reuses every cached slot of category $c$ (over 10k for city names in our corpus) and needs only one respelling call (Sec.~\ref{sec:respell}), with no new speech or slot proposals.
\begin{table*}[t]
  \centering
  \caption{Text-distribution and diversity of 10k outputs per
  generator vs.\ length-matched real CS windows; the real-CS row is
  the reference for all columns and $\dagger$ marks CS-span-masked
  metrics. Func.\ sw.\ = function-word switching (\%).}
  \label{tab:text}
  \small
  \setlength{\tabcolsep}{4pt}
  \begin{tabular}{@{}l ccccccc@{}}
    \toprule
    Method & PPL gap$^\dagger\downarrow$ & CMI gap$\downarrow$
    & MAUVE$\uparrow$ & C2ST$^\dagger$ & Func. sw.
    & Vendi & $\Delta$Self-BLEU$^\dagger$\\
    \midrule
    Real CS speech & 0 & 0 & 1.00 & 0.50 & 0.3 & 4.43 & 0\\
    LLM from scratch & 0.68 & 1.22 & 0.03 & 0.98 & 3.7 & 3.69 & +10.7\\
    Transformation~\cite{liu2026lens}
    & 0.17 & 7.49 & 0.21 & 0.77 & 26.0 & 4.17 & +7.6\\
    RCS & \textbf{0.07} & \textbf{1.13} & \textbf{0.36}
    & \textbf{0.73} & \textbf{0.2} & \textbf{4.39} & \textbf{$-$0.9}\\
    \bottomrule
  \end{tabular}
\end{table*}
\subsection{Phonological spelling augmentation}
\label{sec:respell}
For each catalog value $e$, Gemini~2.5 Flash generates once, offline, a small set of Vietnamese-adapted spellings $\cS(e)=\{s_1,\dots,s_{K_e}\}$ by reasoning from English orthography and pronunciation towards Vietnamese graphemes, optionally with tone marks (e.g., \emph{server}$\rightarrow$\vi{\{sơ-vơ, xơ-vơ\}}, \emph{API}$\rightarrow$\vi{\{ây-pi-ai\}}). Because respellings live at catalog level, their cost is independent of corpus size. Respelling is decided per utterance: with probability $\rho$ the synthesis string $z$ is $y$ with every term $e_k$ replaced by a sampled $s_k\sim\mathcal{U}(\cS(e_k))$, and otherwise $z=y$; the label $y$ always keeps the canonical form. The respelling rate $\rho$ is thus the fraction of synthetic utterances rendered from respelled text ($\rho=0$: canonical TTS; $\rho=1$: always respelled).

Each $z$ is then synthesized end-to-end with the multilingual zero-shot TTS model OmniVoice~\cite{zhu2026omnivoice}, using the source recording as speaker prompt. A lightweight ASR check discards utterances whose recognition deviates from $z$ by more than a threshold, removing gross synthesis failures. The retained pairs form $\cD_{\rm syn}$.

\section{Experimental setup}
\label{sec:setup}
\begin{table}[t]
  \centering
  \caption{Training and evaluation data. EN\%: share of English
  words; CMI: mean code-mixing index; Terms: distinct CS terms.}
  \label{tab:data}
  \small
  \setlength{\tabcolsep}{3.8pt}
  \begin{tabular}{@{}l r r r r r@{}}
    \toprule
    Set & Utts & Hrs & EN\% & CMI & Terms\\
    \midrule
    \multicolumn{6}{@{}l}{\textit{Training}}\\
    GigaSpeech2 VN & 200k & 161.0 & -- & -- & --\\
    RCS & 83.2k & 55.6 & 15.5 & 18.6 & 3{,}718\\
    \midrule
    \multicolumn{6}{@{}l}{\textit{Evaluation}}\\
    VietnamTest & 66.6k & 51.5 & -- & -- & --\\
    GigaSpeech2 CS & 1.5k & 3.9 & 5.5 & 6.0 & 1{,}217\\
    ViMedCSS & 1.6k & 3.4 & 5.9 & 7.1 & 942\\
    \bottomrule
  \end{tabular}
\end{table}

\subsection{Data and text evaluation}
Real training speech comprises 200k Vietnamese utterances from
GigaSpeech2~\cite{yang2025gigaspeech2}. The RCS catalog contains
67 categories and 3,718 terms. From 50.6k source transcripts,
we sample 83.2k synthetic utterances. Unless stated otherwise,
RCS experiments use all 83.2k utterances with $\rho=0.2$.
Evaluation uses GigaSpeech2 CS, the code-switching subset of the
GigaSpeech2 dev set, medical-domain ViMedCSS~\cite{nguyen2026vimedcss},
and the internal monolingual Vietnamese corpus VietnamTest
(Table~\ref{tab:data});
ViMedCSS terms are outside the augmentation catalog.


\begin{table*}[t]
  \centering
  \caption{ASR results (\%). All RCS rows use mixed training;
  Conformer RCS rows use $\rho=0.2$. LM: Vietnamese-only (mono),
  or mono plus RCS transcripts (+RCS). Lex.: base, or base plus
  respelled pronunciations (+resp.). Conformer rows with +RCS LM
  change the LM and lexicon jointly.}
  \label{tab:main}
  \small
  \setlength{\tabcolsep}{3.4pt}
  \begin{tabular}{@{}l ccc ccc cc@{}}
    \toprule
    & \multicolumn{3}{c}{GigaSpeech2 CS}
    & \multicolumn{3}{c}{ViMedCSS}
    & \multicolumn{2}{c}{VietnamTest}\\
    \cmidrule(lr){2-4}\cmidrule(lr){5-7}\cmidrule(l){8-9}
    Training / decoder
    & WER$\downarrow$ & CER$\downarrow$ & TR$\uparrow$
    & WER$\downarrow$ & CER$\downarrow$ & TR$\uparrow$
    & WER$\downarrow$ & CER$\downarrow$\\
    \midrule
    \multicolumn{9}{@{}l}{\textit{Whisper (fine-tuned)}}\\
    Mono (200k)
    & 20.76 & 10.72 & 29.24 & 33.15 & 22.70 & 11.76 & 25.80 & 15.62\\
    +RCS 83.2k, $\rho=0$
    & 18.42 & 9.41 & 39.18 & 32.57 & 21.91 & 14.62 & 23.54 & 14.08\\
    +RCS 83.2k, $\rho=1$
    & 19.03 & 9.69 & 33.21 & 33.08 & 22.24 & 11.93 & 24.06 & 14.39\\
    +RCS 83.2k, $\rho=0.2$
    & \textbf{16.10} & \textbf{7.90} & \textbf{45.30}
    & \textbf{31.95} & \textbf{21.24} & \textbf{17.69}
    & \textbf{22.70} & \textbf{13.50}\\
    \midrule
    \multicolumn{9}{@{}l}{\textit{Conformer-CTC + 4-gram LM (from scratch)}}\\
    Mono / mono LM, base lex.
    & 26.00 & 15.60 & 5.36 & 40.60 & 27.88 & 6.60 & 28.91 & 18.79\\
    Mono / +RCS LM, +resp. lex.
    & 23.32 & 15.26 & 10.12 & 40.50 & 26.59 & 8.10 & 28.81 & 17.89\\
    +RCS / mono LM, base lex.
    & 23.30 & 12.70 & 15.20 & 39.20 & 25.67 & 10.56 & 24.71 & 15.40\\
    +RCS / +RCS LM, +resp. lex.
    & \textbf{20.20} & \textbf{12.36} & \textbf{21.50}
    & \textbf{37.12} & \textbf{24.17} & \textbf{15.69}
    & \textbf{24.02} & \textbf{14.84}\\
    \bottomrule
  \end{tabular}
\end{table*}

For text evaluation, we compare RCS against LLM generation
from scratch and a transformation prompt
following~\cite{liu2026lens}, with 10k outputs per method.
RCS utterances are shorter on average than GigaSpeech2 CS
(2.4\,s versus 9.4\,s; 9.1 versus 37.8 words), and since
EN\% and CMI are normalized by utterance length, their values
in Table~\ref{tab:data} are not directly comparable; text
comparisons therefore use length-matched real CS windows. We report PPL gap, CMI gap,
MAUVE~\cite{pillutla2021mauve}, classifier two-sample test
(C2ST)~\cite{lopezpaz2017c2st}, function-word switching,
Vendi~\cite{friedman2023vendi}, and Self-BLEU.

\subsection{ASR systems and evaluation}
\textbf{Whisper}~\cite{radford2023whisper} denotes
Whisper-large-v3-turbo, fully fine-tuned with AdamW.
\textbf{Conformer-CTC}~\cite{gulati2020conformer,graves2006ctc}
is trained from scratch on Vietnamese onset--rime--tone units
and decoded with a 4-gram word LM and a pronunciation lexicon;
the LM uses Vietnamese-only text (mono), or additionally
the 83.2k RCS transcripts (+RCS), and +resp.\ adds
Vietnamese-adapted lexicon entries. Checkpoints are selected
by WER on the GigaSpeech2 Vietnamese dev set, disjoint from
all test sets.

We report WER, CER, and CS term recall (TR).
Each utterance has a reference term list; a term is recovered
only if its complete form appears in that utterance's
hypothesis. TR measures recovery of reference terms and is
interpreted alongside WER, which also penalizes insertions.
We compare real-only training (Mono), sequential adaptation
on synthetic speech (Mono$\rightarrow$CS), and mixed
training on $\cD_{\rm mono}\cup\cD_{\rm syn}$ (Mix).
Table~\ref{tab:main} uses Mix for all RCS configurations.

\section{Results and discussion}
\label{sec:results}
\subsection{Text distribution and diversity}
RCS is closest to the real-CS reference across the reported
metrics (Table~\ref{tab:text}): its masked PPL gap is 0.07
versus 0.17 for transformation and 0.68 for generation from
scratch, MAUVE is 0.36 versus 0.21 and 0.03, and its
function-word switching rate and Vendi score approach the
reference (0.2\% vs.\ 0.3\%; 4.39 vs.\ 4.43). The higher CMI
in Table~\ref{tab:data} reflects shorter utterances rather
than denser switching, as length-matched windows give 17.9.

C2ST remains above chance at 0.73, and these automatic
metrics do not directly establish perceived naturalness.

\subsection{ASR results across models and domains}
On GigaSpeech2 CS, RCS with $\rho=0.2$ reduces Whisper WER
from 20.76\% to 16.10\% (22.4\% relative) and increases TR
from 29.24\% to 45.30\% (+16.06 percentage points).
Canonical-only RCS already improves WER to 18.42\% and TR
to 39.18\%; adding the mixed-pronunciation configuration
yields a further 2.32-point WER reduction and 6.12-point
TR increase. For Conformer-CTC with the mono LM and base lexicon held fixed, WER
decreases from 26.00\% to 23.30\% and TR increases from 5.36\% to 15.20\%. Thus improvements
are also observed with a recognizer trained from scratch
and without decoder augmentation.

On ViMedCSS, Whisper TR rises from 11.76\% to 17.69\%
(+5.93 points), while WER decreases from 33.15\% to
31.95\% (3.6\% relative); Conformer-CTC TR rises from
6.60\% to 15.69\% and WER decreases from 40.60\% to
37.12\%. This supports transfer beyond the augmentation
catalog within the evaluated setting.
Mixed training also improves Vietnamese-only recognition:
VietnamTest WER decreases from 25.80\% to 22.70\%.

\subsection{Training strategy and respelling}
RCS is applied through mixed training, which avoids the
monolingual degradation observed with sequential adaptation.
Relative to no augmentation, adding RCS at $\rho=0.2$ improves
both WER and CS term recall, and moderate respelling
outperforms both canonical-only and always-respelled
synthesis.

\subsection{Decoder augmentation and data scale}
Jointly adding RCS LM text and respelled lexicon entries
further improves Conformer-CTC TR from 15.20\% to 21.50\%
and WER from 23.30\% to 20.20\% on GigaSpeech2 CS.
Applied to the Mono acoustic model, the same decoder update
yields only 10.12\% TR, well below acoustic augmentation alone
(15.20\%), and the two effects are roughly additive.
The decoder comparison changes the LM and lexicon jointly
and does not isolate the two components.
For Whisper at $\rho=0.2$, TR rises from 29.24\%
without synthetic data to 39.00\%, 42.60\%, and 45.30\%
with 10k, 50k, and 83.2k utterances, respectively.
It continues to improve over the evaluated range
(Fig.~\ref{fig:scale}).
The 10k configuration has numerically similar TR to
83.2k canonical-only utterances (39.00\% versus 39.18\%).

\begin{figure}[t]
  \centering
  \begin{tikzpicture}
  \begin{axis}[
      width=0.98\columnwidth,
      height=3.15cm,
      xlabel={Synthetic utterances},
      ylabel={CS term recall (\%)},
      xtick={0,10,50,83.2},
      xticklabels={0,10k,50k,83.2k},
      ymin=26, ymax=48,
      ymajorgrids,
      grid style={gray!25},
      tick label style={font=\small},
      label style={font=\small},
      nodes near coords,
      every node near coord/.append style={font=\small,anchor=south},
  ]
  \addplot+[mark=*, thick] coordinates {
      (0,29.24) (10,39.00) (50,42.60) (83.2,45.30)
  };
  \end{axis}
  \end{tikzpicture}
  \vspace{-0.7em}
  \caption{CS term recall of Whisper on GigaSpeech2 CS versus the number of
synthetic utterances ($\rho{=}0.2$). All synthetic sets are drawn from the
same monolingual source corpus.}
  \label{fig:scale}
  \vspace{-0.9em}
\end{figure}

The evaluation covers one language pair and one TTS model;
matched ASR comparisons against alternative text generators
are not included.

\section{Conclusion}
RCS constructs Vietnamese--English ASR training data by
combining catalog-grounded slot editing of real transcripts
with Vietnamese-adapted respelling. Mixed training improves
the reported recognition metrics for Whisper and
Conformer-CTC, including results on a Vietnamese-only test
set. These results support RCS as a practical augmentation
approach within
the evaluated setting; broader downstream comparisons
and perceptual validation remain open.

\bibliographystyle{IEEEbib}
\bibliography{refs}
\end{document}